%! AP Statistics — Unit 4, Lesson 4.10 — Homework
\documentclass[10pt]{article}
\usepackage{apstats-article}
\usepackage{apstats-boxes}
\usepackage{grading}
\gradingassignment{apstats-u04-l10-hw}
\gradingcourse{advanced_statistics}

\begin{document}

\pageheader{Unit 4, Lesson 4.10}{Homework}
\namedateperiod

\vspace{0.1in}

\begin{practicebox}
\small
\textbf{Directions:} For each scenario, (a) check BINS and state whether the situation
is a binomial setting, and (b) if yes, compute the requested probability showing the
formula setup.

\vspace{10pt}
\textbf{1.} A basketball player makes 72\% of her free throws. She shoots 10 free
throws in a game. Let $X =$ the number she makes.
\begin{enumerate}[label=(\alph*), leftmargin=2em, itemsep=8pt]
  \item Check BINS and state the distribution (or explain which condition fails).\\
        \gradedhide{\writeline\gradedright{\raisebox{-13pt}{\gradedblank[id=hw1a, type=rubric, points=3, display={B, I, N, S all met (fixed $n = 10$, same $p = 0.72$, independent shots): $X \sim B(10, 0.72)$.}]{3.3in}{\gradedlinesheight{3}-5pt}}}\\[1pt]\writeline\\[1pt]\writeline}
  \item Find $P(X = 8)$. Show the formula and the numerical answer.\\
        \writeline\\[1pt]\writeline\gradedhide{\hrulefill}\gradedright{\raisebox{-2pt}{\gradedblank[id=hw1b, type=short, points=2, answer=0.2548, check=numeric, tolerance=0.002]{1.3in}{0.2in}}}
\end{enumerate}

\vspace{8pt}
\textbf{2.} A survey asks 200 randomly selected adults (from a city of 500{,}000)
whether they support a new recycling program. Historically 60\% support such
programs. Let $X =$ the number who support the program.
\begin{enumerate}[label=(\alph*), leftmargin=2em, itemsep=8pt]
  \item Check BINS. Is the $10\%$ condition relevant here? Explain.\\
        \gradedhide{\writeline\gradedright{\raisebox{-5pt}{\gradedblank[id=hw2a, type=rubric, points=2, display={BINS met; yes, the 10\% condition matters (sampling without replacement): $200 \le 50{,}000$.}]{5.2in}{\gradedlinesheight{2}-5pt}}}\\[1pt]\writeline}
  \item Find $P(X = 120)$ using \texttt{binompdf}. (No hand calculation required ---
        just set up the calculator notation and record the result.)\\
        \writeline\gradedhide{\hrulefill}\gradedright{\raisebox{-2pt}{\gradedblank[id=hw2b, type=short, points=1, answer=0.0553, check=numeric, tolerance=0.001]{1.3in}{0.2in}}}
\end{enumerate}

\vspace{8pt}
\textbf{3.} A class has 30 students. The teacher randomly calls on 6 students without
replacement. 10 of the 30 students did not complete the homework.
Let $X =$ the number called on who did not complete the homework.
\begin{enumerate}[label=(\alph*), leftmargin=2em, itemsep=8pt]
  \item Is this a binomial setting? Which condition(s), if any, are violated?\\
        \gradedhide{\writeline\gradedright{\raisebox{-5pt}{\gradedblank[id=hw3a, type=rubric, points=2, display={Not binomial: drawing without replacement, so $p$ changes from call to call (not independent).}]{6.1in}{\gradedlinesheight{2}-5pt}}}\\[1pt]\writeline}
  \item Check the $10\%$ condition: is $n \le 0.10 \times N$? Show the calculation.\\
        \gradedhide{\writeline\gradedright{\raisebox{-1pt}{\gradedblank[id=hw3b, type=rubric, points=2, display={$0.10 \times 30 = 3 < 6$: condition fails.}]{3.3in}{0.22in}}}}
\end{enumerate}

\vspace{8pt}
\textbf{4.} Evaluate $\binom{6}{2}$ by hand. Then explain in one sentence what this
number represents in the context of a binomial experiment with $n = 6$ trials where
you are interested in exactly 2 successes.
\begin{enumerate}[label=(\alph*), leftmargin=2em, itemsep=8pt]
  \item $\binom{6}{2} = $ \makebox[3cm][l]{\hspace{1.6cm}\smash{\raisebox{-4pt}{\gradedblank[id=hw4a, type=short, points=1, answer=15, check=numeric, tolerance=0]{1.3cm}{0.2in}}}}
  \item Interpretation: \gradedhide{\writeline\gradedright{\raisebox{-5pt}{\gradedblank[id=hw4b, type=rubric, points=2, display={The number of ways (orderings) to get exactly 2 successes in 6 trials.}]{4.8in}{\gradedlinesheight{2}-5pt}}}\\[1pt]\writeline}
\end{enumerate}

\vspace{8pt}
\textbf{5.} (AP-Style Free Response) \textit{A door-to-door salesperson visits
8 houses. Based on past experience, the probability of making a sale at any given
house is 0.15. Assume houses are independent.}
\begin{enumerate}[label=(\alph*), leftmargin=2em, itemsep=8pt]
  \item Define the random variable $X$ and state its distribution. Check BINS
        briefly.\\
        \gradedhide{\writeline\gradedright{\raisebox{-5pt}{\gradedblank[id=hw5a, type=rubric, points=2, display={$X$ = number of sales in 8 houses; BINS met; $X \sim B(8, 0.15)$.}]{3.6in}{\gradedlinesheight{2}-5pt}}}\\[1pt]\writeline}
  \item Calculate $P(X = 1)$. Show formula setup.\\
        \writeline\\[1pt]\writeline\gradedhide{\hrulefill}\gradedright{\raisebox{-2pt}{\gradedblank[id=hw5b, type=short, points=2, answer=0.3847, check=numeric, tolerance=0.002]{1.3in}{0.2in}}}
  \item Calculate $P(X = 0)$. Interpret in context.\\
        \writeline\gradedhide{\hrulefill}\gradedright{\raisebox{-2pt}{\gradedblank[id=hw5c1, type=short, points=1, answer=0.2725, check=numeric, tolerance=0.002]{1.3in}{0.2in}}}\\[1pt]\gradedhide{\writeline\gradedright{\raisebox{-1pt}{\gradedblank[id=hw5c2, type=rubric, points=2, display={About a 27.3\% chance the salesperson makes no sales in the 8 houses.}]{6.1in}{0.22in}}}}
\end{enumerate}

\vspace{8pt}
\textbf{6. Extension:} Build the complete probability distribution table for
$X \sim B(5, 0.4)$ by computing $P(X = k)$ for $k = 0, 1, 2, 3, 4, 5$.
Verify that the probabilities sum to 1.

\vspace{6pt}
\renewcommand{\arraystretch}{1.8}
\begin{tabularx}{\linewidth}{@{} >{\bfseries}p{1.2cm} X X X X X X @{}}
$k$ & 0 & 1 & 2 & 3 & 4 & 5 \\
\hline
$P(X=k)$ & \blank{1.5cm} & \blank{1.5cm} & \blank{1.5cm} & \blank{1.5cm} & \blank{1.5cm} & \blank{1.5cm} \\
\end{tabularx}

\vspace{6pt}
Sum $= $ \blank{2cm} (should equal 1)

\end{practicebox}

\end{document}
